\exercisesfor{5}

第 5 节的约定除非另有说明均有效。

\begin{enumerate}
\def\labelenumi{\arabic{enumi}.}
\item
  设 g 为 2 维非交换可解李代数。证明 g 的 Killing 形式非零，g 上的每个不变双线性形式都退化，且 g 的每个导子都是内导子。
\item
  \begin{enumerate}
  \def\labelenumii{(\alph{enumii})}
  \tightlist
  \item
    证明在 \(\mathbf{R}\) 上由乘法表 \([x, y] = z\)、\([x, z] = -y\)、\([y, z] = 0\) 定义的 3 维可解李代数中，不存在维数为 3、2、1、0 的递降理想序列。
  \end{enumerate}
\end{enumerate}

\begin{enumerate}
\def\labelenumi{(\alph{enumi})}
\setcounter{enumi}{1}
\tightlist
\item
  证明在非交换的 2 维可解李代数 g 中存在维数为 2、1、0 的理想序列，但 g 并非幂零。
\end{enumerate}

\begin{enumerate}
\def\labelenumi{\arabic{enumi}.}
\setcounter{enumi}{2}
\item
  设 \(\mathfrak{g}\) 为可解李代数，使得条件 \(x\in \mathfrak{g},y\in \mathfrak{g}\)、\([[x,y],y] = 0\) 蕴含 \([x,y] = 0\)。证明 \(\mathfrak{g}\) 是交换的。（设 \(k\) 为使 \(\mathcal{D}^{k - 1}\mathfrak{g}\neq \{0\}\)、\(\mathcal{D}^k\mathfrak{g} = \{0\}\) 的最大整数。假设 \(k\geqslant 2\)，先证明 \([\mathcal{D}^{k - 2}\mathfrak{g},\mathcal{D}^{k - 1}\mathfrak{g}] = \{0\}\)，再证明 \([\mathcal{D}^{k - 2}\mathfrak{g},\mathcal{D}^{k - 2}\mathfrak{g}] = \{0\}\)，从而得出矛盾。）
\item
  证明 \(\mathfrak{st}(n,\mathbf{K})\) 的中心为零，而 \(\mathfrak{n}(n,\mathbf{K})\) 的中心维数为 1。
\item
  设 \(\mathfrak{g}\) 为李代数，\((\mathcal{D}^0\mathfrak{g},\mathcal{D}^1\mathfrak{g},\ldots ,\mathcal{D}^n\mathfrak{g})\) 为 \(\mathfrak{g}\) 的导代数序列（ \(n\geqslant 0,\mathcal{D}^{n - 1}\mathfrak{g}\neq \mathcal{D}^n\mathfrak{g}\) ）。则 \(\dim \mathcal{D}^i\mathfrak{g} / \mathcal{D}^{i + 1}\mathfrak{g}\geqslant 2^{i - 1} + 1\) 对 \(1\leqslant i\leqslant n - 2\) 成立。（对 \(\mathcal{D}^n\mathfrak{g}\) 取商，归结为 \(\mathfrak{g}\) 可解的情形。然后利用 \(\mathcal{D}_{\mathfrak{g}}\) 幂零这一事实以及第 4 节练习 7 (c)。）
\item
  \begin{enumerate}
  \def\labelenumii{(\alph{enumii})}
  \tightlist
  \item
    验证以下乘法表定义了一个 5 维可解李代数 g：
  \end{enumerate}
\end{enumerate}

\[
\begin{array}{r} \left[ x _ {1}, x _ {2} \right] = x _ {5}, \qquad \left[ x _ {1}, x _ {3} \right] = x _ {3}, \qquad \left[ x _ {2}, x _ {4} \right] = x _ {4} \\ \left[ x _ {1}, x _ {4} \right] = \left[ x _ {2}, x _ {3} \right] = \left[ x _ {3}, x _ {4} \right] = \left[ x _ {5}, \mathfrak {g} \right] = 0. \end{array}
\]

\begin{enumerate}
\def\labelenumi{(\alph{enumi})}
\setcounter{enumi}{1}
\item
  证明 \(\mathfrak{g}\) 关于 Killing 形式的正交子空间是 \(\mathcal{D}\mathfrak{g} = \mathrm{K}x_3 + \mathrm{K}x_4 + \mathrm{K}x_5\)。由此推出 \(\mathcal{D}\mathfrak{g}\) 是 \(\mathfrak{g}\) 的最大幂零理想。
\item
  证明不存在 \(\mathcal{D}\mathfrak{g}\) 在 \(\mathfrak{g}\) 中的补子代数。由此推出 \(\mathfrak{g}\) 不是交换代数与一个幂零理想的半直积。（证明该幂零理想必然就是 \(\mathcal{D}\mathfrak{g}\)。）
\end{enumerate}

\begin{enumerate}
\def\labelenumi{\arabic{enumi}.}
\setcounter{enumi}{6}
\item
  设 \(\mathfrak{g}\) 为以 \((x,y,z)\) 为基、满足 \([x,y] = z\)、\([x,z] = y\)、\([y,z] = 0\) 的 3 维可解李代数。证明将 \(x\) 映为 \(-x, y\) 映为 \(-z, z\) 映为 \(y\) 的线性映射是 \(\mathfrak{g}\) 的一个 4 阶自同构。将此结果与第 4 节练习 21 (c) 相比较。
\item
  \begin{enumerate}
  \def\labelenumii{(\alph{enumii})}
  \tightlist
  \item
    设 \(\mathfrak{g}_0\) 为 3 维实可解李代数，使得 \(\mathcal{D}\mathfrak{g}_0\) 交换且 2 维。设 \(\mathfrak{g}\) 为由 \(\mathfrak{g}_0\) 通过将基域从 \(\mathbf{R}\) 扩张到 \(\mathbf{C}\) 而得到的代数。对 \(x \in \mathfrak{g}\)，令 \(u_x\) 为 \(\mathrm{ad}_{\mathfrak{g}}x\) 在 \(\mathcal{D}\mathfrak{g}\) 上的限制。证明 \(u_x\) 的特征值或者具有相同的绝对值，或者在 \(\mathbf{R}\) 上线性相关。（我们有 \(x = \lambda y + z\)，其中 \(z \in \mathcal{D}\mathfrak{g}\)、\(y \in \mathfrak{g}_0\)、\(\lambda \in \mathbf{C}\)，因而 \(u_x = \lambda u_y\)。而 \(u_y\) 是 \(\mathbf{C}\) 线性地扩张 \(\mathcal{D}\mathfrak{g}\) 的一个 \(\mathbf{R}\) 线性自同态所得的映射，后者定义在 \(\mathcal{D}\mathfrak{g}_{0}\) 上。）
  \end{enumerate}
\end{enumerate}

\begin{enumerate}
\def\labelenumi{(\alph{enumi})}
\setcounter{enumi}{1}
\item
  证明存在一个 3 维复可解李代数 g，其 \(\mathcal{D}_{\mathfrak{g}}\) 交换且 2 维，并且存在 g 的元素 \(x\)，使得 \(\mathrm{ad}_{\mathfrak{g}}x\) 在 \(\mathcal{D}_{\mathfrak{g}}\) 上的限制的特征值既不具有相同的绝对值，也不在 R 上线性相关。（将 g 构造为 1 维代数与 2 维交换代数的半直积。）
\item
  证明 (b) 中构造的代数不能由某个实李代数通过将基域从 \(\mathbf{R}\) 扩张到 \(\mathbf{C}\) 而得到。
\end{enumerate}

\begin{enumerate}
\def\labelenumi{\arabic{enumi}.}
\setcounter{enumi}{8}
\tightlist
\item
  设 \(\mathfrak{g}\) 为李代数，\(\mathfrak{r}\) 为其根基，\(\mathfrak{n}\) 为其最大幂零理想，D 为 \(\mathfrak{g}\) 的一个导子。证明 \(\mathcal{D}\mathfrak{g} \cap \mathfrak{r} \subset \mathfrak{n}\)。（设 \(\mathfrak{d}\) 为 \(\mathfrak{g}\) 的导子组成的李代数，\(\mathfrak{r}'\) 为其根基。若 \(x \in \mathfrak{g}\) 使得 \(\mathrm{D}x \in \mathfrak{r}\)，则
\end{enumerate}

\[
\operatorname{ad} (\mathrm{D} x) = [ \mathrm{D}, \operatorname{ad} x ] \in \mathscr {D} \mathfrak {d} \cap \mathfrak {r} ^ {\prime}
\]

（命题 5 的推论 2），因而 \(\operatorname{ad}(\mathbf{D}x)\) 幂零（定理 1），从而 \(\mathbf{D}x \in \mathfrak{n}\)。）

\begin{enumerate}
\def\labelenumi{\arabic{enumi}.}
\setcounter{enumi}{9}
\item
  设 g 为李代数，r 为其根基，a 为 g 的一个次不变子代数。证明 a 的根基是 \(a \cap r\)。（若干次应用命题 5 的推论 3。）
\item
  设 \(\mathfrak{g}\) 为李代数，\(\rho\) 为 \(\mathfrak{g}\) 的有限维表示，\(\mathbb{E}\) 为由 1 与 \(\rho(\mathfrak{g})\) 生成的结合自同态代数。证明理想 \(n\)（\(\rho\) 的最大幂零性理想）等于集合 \(n'\)，它由 \(x \in \mathfrak{g}\)（满足 \(\operatorname{Tr}(\rho(x)u) = 0\)，对一切 \(u \in E\)）组成。（为证明 \(n' \subset n\)，证明 \(n'\) 是一个理想，且对一切 \(x \in n'\)，\(\rho(x)\) 的半单分量为零，这通过注意 \(\operatorname{Tr}((\rho(x))^n) = 0\) 对每个整数 \(n > 0\) 成立得到。）
\item
  假设 K 代数闭。设 g 为幂零李 K-代数。设 \(\rho\) 为 g 在向量空间 V 上的一个有限维表示。对 g 上的每个线性形式 \(\lambda\)，令 \(V^{\lambda}\) 为 V 中由满足如下条件的 \(\xi \in V\) 组成的向量子空间：对一切 \(x \in g\)，当 n 充分大时 \((\rho(x) - \lambda(x)\mathrm{I})^{n}\xi = 0\)。
\end{enumerate}

\begin{enumerate}
\def\labelenumi{(\alph{enumi})}
\item
  诸 \(V^{\lambda}\) 在 \(\rho(\mathfrak{g})\) 下稳定，且诸 \(V^{\lambda}\) 之和是直和。（使用第 4 节练习 22。）
\item
  我们有 \(\mathrm{V} = \sum \mathrm{V}^{\lambda}\)。（若每个 \(\rho(x)\) 都只有一个特征值，则由定理 1 的推论 2 得 \(\mathrm{V} = \mathrm{V}^{\lambda_0}\)。若 \(\rho(x_0)\) 至少有两个不同的特征值，则 \(\mathrm{V}\)）
\end{enumerate}

是两个在 \(\rho(\mathfrak{g})\) 下稳定的非平凡子空间的直和。然后对 V 的维数作归纳。）

\begin{enumerate}
\def\labelenumi{\arabic{enumi}.}
\setcounter{enumi}{12}
\item
  设 \(\mathfrak{g}\) 为李代数，\(\mathfrak{D}\) 为 \(\mathfrak{g}\) 的导子组成的李代数。\(\mathfrak{g}\) 特征幂零的充分必要条件是 \(\mathfrak{D}\) 幂零且 \(\dim \mathfrak{g} > 1\)。（为验证条件充分，把 \(\mathfrak{g}\) 写成子空间 \(\mathfrak{g}^{\lambda}\) 之和，办法是把练习 12 应用于 \(\mathfrak{D}\) 的恒等映射。证明 \([\mathfrak{g}^{\lambda}, \mathfrak{g}^{\mu}] \subset \mathfrak{g}^{\lambda + \mu}\)，且每个 \(\mathfrak{g}^{\lambda}\) 都是 \(\mathfrak{g}\) 的理想。由此推出 \(\mathfrak{g}^{\lambda}\) 当 \(\lambda \neq 0\) 时是交换的。再利用 \(\mathfrak{D}\) 幂零这一事实，证明 \(\mathfrak{g} = \mathfrak{g}^{0}\)（若 \(\dim \mathfrak{g} > 1\)）；否则 \(\mathfrak{g} = \mathfrak{g}^{0} \times \mathfrak{h}\)，其中 \(\mathfrak{h}\) 交换。先证明 \(\dim \mathfrak{h} \leqslant 1\)。若 \(\dim \mathfrak{h} = 1\)，注意将存在 \(\mathfrak{g}\) 的导子 D，使得 D( \(\mathfrak{g}^{0}\) ) = \{0\} 且 D( \(\mathfrak{h}\) ) 包含在 \(\mathfrak{g}^{0}\) 的中心中（ \(\S 4\)，练习 8 ( \(a\) )）。）
\item
  设 V 为 K 上有限维向量空间，z 为 V 的自同态。我们采用第 3 节练习 3 的记号 \(z_{q}^{p}\)。V 的自同态 \(z'\) 若对一切 p 和 q 使得 \(z_{q}^{p}\) 的每个零点都是 \(z_{q}^{\prime p}\) 的零点，则称为 z 的一个复本。证明若 \(\operatorname{Tr}(zz') = 0\) 对每个复本 \(z'\)（\(z\) 的复本）成立，则 z 幂零。（利用引理 3 的证明。在该证明的记号下，特别证明 t 是 z 的复本。）
\item
  设 K 为特征 2 的域。幂零李代数 \(\mathfrak{sl}(2,\mathbf{K})\) 在 \(K^{2}\) 上的恒等表示定义了 \(\mathfrak{sl}(2,\mathbf{K})\) 与 \(K^{2}\) 的一个半直积 h。证明 h 可解，但 Dh 不幂零。由此推出 h 不容许任何以三角矩阵给出的忠实线性表示。再证明练习 5 的结论在此为假。
\item
  设 g 为任意域 K 上的李代数，A 为 K 上交换结合代数，而 \(g' = g \otimes_{K} A\)，它可视为 K 上的李代数。
\end{enumerate}

\begin{enumerate}
\def\labelenumi{(\alph{enumi})}
\item
  若 D 是 A 的导子，证明存在唯一的导子 \(D'\)（\(g'\) 的导子），使得 \(\mathrm{D}'(x \otimes a) = x \otimes \mathrm{D}a\) 对 \(x \in g, a \in A\) 成立。
\item
  设 \(p\) 为素数，\(G\) 为 \(p\) 阶循环群，\(s\) 为 \(G\) 的一个生成元。假设 \(K\) 的特征为 \(p\)，以下设 \(A\) 为 \(G\) 在 \(K\) 上的代数。证明存在唯一的导子 \(D\)（\(A\) 的导子），使得 \(D(s^k) = ks^{k-1}\) 对一切 \(k \in \mathbf{Z}\) 成立。证明取 \(K\) 中系数时诸 \(x \otimes (s - 1)^k\)（ \(k = 1, 2, \ldots, p - 1, x \in \mathfrak{g}\) ）的线性组合构成一个可解理想 \(\tau\)（\(g'\) 的理想），且 \(g'/\tau\) 同构于 \(\mathfrak{g}\)。
\item
  取 \(\mathfrak{g}\) 为单代数（参见~第 6 节，第 2 号，定义 2）。则 \(\mathfrak{r}\) 是 \(\mathfrak{g}'\) 的根基，但不是特征理想。（注意 \(\mathrm{D}(x\otimes (s - 1)) = 1\otimes x.\) ）
\end{enumerate}

\begin{enumerate}
\def\labelenumi{\arabic{enumi}.}
\setcounter{enumi}{16}
\tightlist
\item
  设 g 为李代数。假设对每个 K 上有限维单 g-模 M，诸 \(x_{M}\) 两两可交换。证明 g 可解。（注意 Dg 包含在幂零根基中，因而可解。）
\end{enumerate}

